\label{sec:outlook}
\subsection{What the placement explains}
\begin{enumerate}[leftmargin=*]
\item \textbf{Why the mean is hard and the densities are easy.} The mean is a sum over walls of a density at the kink
times a conditional transport on the wall (Theorem~\ref{thm:kink}). Densities are concentration-function
coefficients and the Gaussian surrogate gets them right; the conditional transports are rare-event-relative objects
(conditioning on a measure-zero wall), and every chain's error lives there.
\item \textbf{Why the third-cumulant memory has an identity leg.} It is the label of the kink at which the cumulant was
born (Theorem~\ref{thm:birth}); the memory is the record of all pinnings, transported.
\item \textbf{Why it cannot be compressed.} Kinks have small codegree at He initialisation, so their atoms are orthogonal;
the frame bound (Theorem~\ref{thm:frame}) makes the memory exactly as compressible as its transport leg, and the
complementarity of the neuron algebra with the modular frame of the layer state (Theorem~\ref{thm:complement}) makes
state-respecting (spectral) compressions blind to the hub diagonal up to $2\sqrt{2/n}$. The same small codegree that
makes Sahasrabudhe's nibble concentrate makes our compression fail.
\item \textbf{Why ensemble surrogates fail.} The memory's readout is odd in the last weight matrix
(Proposition~\ref{prop:odd}); free and NNGP descriptions are the symmetric phase.
\item \textbf{Where the Euler--Stein constants come from.} The correction averages the defect over Gaussian balls of
uniform radius (Theorem~\ref{thm:uniform}); a defect of Hermite contact order $d$ decays like $r^{d+1}$, so
$a=1/(d+2)$, and $d$ is the order of the concentration coefficient the estimator gets wrong
(Theorem~\ref{thm:order}, Proposition~\ref{prop:conc}); the three measured constants fit $d=0,0,2$.
\item \textbf{Why depth behaves algebraically.} Lengths are preserved and chords contract like $3\pi/l$
(Theorem~\ref{thm:folding}) because the activation's heat trace is critical at linear order and has a
$\beta^{3/2}$ kink term (Proposition~\ref{prop:heattrace}).
\end{enumerate}

\subsection{What it predicts (falsifiable, theoretical)}
\begin{enumerate}[leftmargin=*]
\item The defect of the exact first-order chain along Eldan's path decays like $r^3=(1+\tau)^{-3/2}$ (order $d=2$),
not like $(1+\tau)^{-1}$; a chain that also corrects the curvature coefficient $f''(0)$ at the kinks has a defect of order
$d\ge4$ and merge constant at most $1/6$.
\item On the official networks, the number of walls crossed by a great-circle arc grows linearly in depth (about
$n\theta_0/\pi$ per layer) while the number of gates that differ between its endpoints grows like $3n\log L$.
\item The best hub-index compression of any old source has the error of the low-rank truncation of its transport leg
$T_{l\leftarrow l'}C$ (Theorem~\ref{thm:frame} with a near-isotropic Hadamard leg); a tier cannot do better than that number,
which can be computed without running any chain.
\item The complementarity constant of the layer state at depth is $\approx2\sqrt{2/n}$ up to the collective mode, so every
state-preserving coarse-graining of a layer onto spectral data retains at most a fraction $\approx8/n$ of the centred
energy of any neuron-diagonal observable (a gate pattern, a hub diagonal).
\end{enumerate}

\subsection{Open problems}
\begin{enumerate}[leftmargin=*]
\item \textbf{Contact order at depth.} Prove Conjecture~\ref{conj:order}: show that the defect of a chain whose first
neglected term is of Hermite order $d$ at every wall has the profile $r^{d+1}$ along the localization, despite the
nonlinear dependence of the downstream pre-activation means on the localization centre.
\item \textbf{A contact-resolved estimator.} The kink-current formula asks for the on-wall transports
$\E[\partial_{h_l}h_L\,e_a\mid z_{l,a}=0]$ for all walls. To first order these are pinning responses
(Lemma~\ref{lem:pinning}); their sum over walls is one directional derivative of the downstream chain per output,
which forward mode computes with an identity leg (the memory) and reverse mode with one adjoint per output. Either way
the cost is that of the memory. Is there an order-two quantity (a pinning response of pinning responses) whose sum
over walls collapses, as the radial fourth-cumulant scalar does for the fourth cumulant?
\item \textbf{The induced Dirac ensemble.} Compute the spectral dimension and spectral variance \cite{BDG} of the
selected Dirac operators $P_xDP_x$ and of the transports $T_{l\leftarrow l'}$ by free multiplicative convolution, and
compare with the measured stable ranks $164,92,65,47,29,23$ of the cocycle at ages $1,2,3,4,6,8$; decide whether the
emergence of the collective eigenvalue of the layer covariance (from $0.9\%$ to $37$--$48\%$ of the energy) is a
Baik--Ben Arous--P\'ech\'e transition in depth, the analogue of the one-cut to two-cut transition of \cite{DG}.
\item \textbf{Bootstrap with kink generators.} Is the island of the moment problem at low degree narrow at large $n$
because of a factorisation analogous to large-$N$ factorisation \cite{KP}? If so, positivity would certify chains
cheaply.
\item \textbf{Klartag in reverse.} Klartag's ellipsoid stops when about $d^2/2$ contacts hold its $d(d+1)/2$ shape
parameters. The marginal value of a source age decays by a factor $0.75$ per age and about eleven ages matter,
roughly $4n$ effective contacts held in memory. Is there a counting argument that fixes this number from the dimension
of the third-cumulant readouts that the mean actually uses?
\end{enumerate}
